\documentclass[10pt,letterpaper]{article}

\usepackage[margin=0.58in,top=0.50in,bottom=0.55in]{geometry}
\usepackage[T1]{fontenc}
\usepackage[utf8]{inputenc}
\usepackage{lmodern}
\usepackage{microtype}
\usepackage{xcolor}
\usepackage{tabularx}
\usepackage{enumitem}
\usepackage{titlesec}
\usepackage{fancyhdr}
\usepackage{lastpage}
\usepackage[hidelinks]{hyperref}
\usepackage{fontawesome5}

\definecolor{linkblue}{HTML}{1E5AA8}
\definecolor{sectionblue}{HTML}{1565C0}

\hypersetup{
  colorlinks=true,
  urlcolor=linkblue,
  linkcolor=black
}

\pagestyle{fancy}
\fancyhf{}
\renewcommand{\headrulewidth}{0pt}
\renewcommand{\footrulewidth}{0pt}
\fancyfoot[C]{\small Page \thepage\ of \pageref{LastPage}\hspace{1.2em}\thepage/\pageref{LastPage}}

\setlength{\parindent}{0pt}
\setlength{\parskip}{0pt}
\setlength{\tabcolsep}{0pt}
\renewcommand{\arraystretch}{1.05}
\setlist[itemize]{leftmargin=1.25em,itemsep=0.10em,topsep=0.15em,parsep=0pt,partopsep=0pt,label=\raisebox{0.15ex}{\scriptsize$\circ$}}

\titleformat{\section}
  {\large\bfseries\color{sectionblue}}
  {}
  {0pt}
  {\textcolor{sectionblue}{\rule[0.45ex]{1.65em}{0.46ex}}\hspace{0.45em}}
  [\vspace{-0.35em}]
\titlespacing*{\section}{0pt}{0.70em}{0.42em}

\newcommand{\cvsection}[1]{\section*{#1}}

\newcommand{\daterow}[3]{%
  \begin{tabularx}{\textwidth}{@{}p{1.18in}X@{}}
    #1 & #2\\
       & #3
  \end{tabularx}\vspace{0.18em}
}

\newcommand{\datedetailrow}[4]{%
  \begin{tabularx}{\textwidth}{@{}p{1.18in}X@{}}
    #1 & \textbf{#2}\\
       & #3\\
       & #4
  \end{tabularx}\vspace{0.18em}
}

\newcommand{\publication}[3]{%
  \begin{tabularx}{\textwidth}{@{}p{1.18in}X@{}}
    \textbf{#1} & #2\\
                & #3
  \end{tabularx}\vspace{0.22em}
}

\newcommand{\service}[2]{%
  \begin{tabularx}{\textwidth}{@{}p{1.18in}X@{}}
    \textbf{#1} & #2
  \end{tabularx}\vspace{0.22em}
}

\newcommand{\compactitem}[1]{\item #1}

\begin{document}

\begin{tabularx}{\textwidth}{@{}Xr@{}}
  {\LARGE\bfseries Sergei Makeev} &
  \begin{minipage}[t]{0.48\textwidth}
    \raggedleft
    \faGlobe\ \href{https://neuralsrg.github.io}{neuralsrg.github.io}\\[-0.05em]
    \faEnvelope\ \href{mailto:neuralsrg@gmail.com}{neuralsrg@gmail.com}\\[-0.05em]
    \faGithub\ \href{https://github.com/neuralsrg}{github.com/neuralsrg}\\[-0.05em]
    \faLinkedin\ \href{https://www.linkedin.com/in/sergei-makeev}{linkedin.com/in/sergei-makeev}
  \end{minipage}
\end{tabularx}

\vspace{0.30em}

\cvsection{Research Interests}
Information retrieval, reinforcement learning, recommender systems, large language models, generative models and scaling laws, and the mathematical foundations of deep learning.

\cvsection{Education}
\datedetailrow{2024 -- 2026}{M.Sc., Applied Mathematics and Computer Science, Moscow State University}{Faculty of Computational Mathematics and Cybernetics. Diploma with Honours.}{Thesis topic: Approximation of softmax activation function in recommendation problem.}
\datedetailrow{2020 -- 2024}{B.Sc., Applied Mathematics and Computer Science, Moscow State University}{Faculty of Computational Mathematics and Cybernetics. Diploma with Honours.}{Thesis topic: Application of discrete mathematics and combinatorics in machine learning.}

\cvsection{Work Experience}
\datedetailrow{2026 -- Present}{Senior Researcher, AI VK}{Part of a research team led by \href{https://www.kaggle.com/dyakonov}{Prof. Dr. Alexander D'yakonov}, formerly world's \#1 rank on Kaggle. Research topics: generative retrieval and reinforcement learning in recommender systems.}{
  \begin{itemize}
    \compactitem{Co-authored \href{https://neuralsrg.github.io/publication/2026-07-01-long-term-optimization}{Long-Term Optimization for Large-Scale Generative Retrieval with Off-Policy REINFORCE}, accepted to the 5th Workshop on End-to-End Customer Journey Optimization at KDD 2026.}
  \end{itemize}
}
\datedetailrow{2024 -- 2026}{Senior Deep Learning Engineer, RecSys R\&D, Yandex}{Worked under the leadership of \href{https://www.linkedin.com/in/kirill-khrylchenko/}{Kirill Khrylchenko}. Research topics: generative models and scaling laws, and linear attention in real-time recommendation models.}{
  \begin{itemize}
    \compactitem{Part of the team that developed \href{https://neuralsrg.github.io/publication/2026-07-05-argus}{Argus}. Argus increased TLT at Yandex Music by +2.26\% and like likelihood by +6.37\%. It was later adapted by our colleagues to nearly every recommendation-based Yandex product (e.g., e-commerce +1.7\% GMV, food-tech services +2.1\% GMV, and others).}
    \compactitem{Worked on quality-complexity trade-offs for recommendation models, co-authoring \href{https://neuralsrg.github.io/publication/2026-06-08-gbla}{Gated Bidirectional Linear Attention for Generative Retrieval}.}
    \compactitem{Worked on pretraining large-scale generative recommender models.}
  \end{itemize}
}
\datedetailrow{2024}{Research Intern, RecSys R\&D, Yandex}{Research topic: applications of graph neural networks in web-scale recommender systems.}{}
\datedetailrow{2021 -- 2024}{Research Assistant, Moscow State University}{Research topics: cognitive dynamics analysis and internal speech recognition under Russian Science Foundation grants \href{https://rscf.ru/project/LXjuaL8yl_MKyMwV-E0AYmYXKZ7dVATTIy5v_6WfDOq3PZJVv6H2gRUOxla_daUJ4bgTN-xp/}{\#20-18-00067} and \href{https://rscf.ru/project/OwNy03bhBYAiNgcRD1Lq9DYlzpziFaDtdi0LjD-lDfhmGp1fntBfNIA3DdjvFBf32QgMbnrg/}{\#20-18-00067-\(\Pi\)}, under the supervision of \href{https://www.researchgate.net/profile/Vartanov-Alexander}{Dr. Alexander Vartanov}.}{}

\cvsection{Publications}
\publication{KDD 2026}{\href{https://neuralsrg.github.io/publication/2026-07-05-argus}{Scaling Recommender Transformers to One Billion Parameters}}{Kirill Khrylchenko, Artem Matveev, \textbf{Sergei Makeev}, Vladimir Baikalov}
\publication{}{\href{https://neuralsrg.github.io/publication/2026-07-01-long-term-optimization}{Long-Term Optimization for Large-Scale Generative Retrieval with Off-Policy REINFORCE}, \textit{5th Workshop on End-to-End Customer Journey Optimization}}{Artem Matveev, \textbf{Sergei Makeev}, Aleksei Krasilnikov, Vladimir Baikalov, Sergei Liamaev, Kirill Khrylchenko}
\publication{SIGIR 2026}{\href{https://neuralsrg.github.io/publication/2026-06-08-gbla}{Gated Bidirectional Linear Attention for Generative Retrieval}}{Artem Matveev, Vladislav Tytskiy, \textbf{Sergei Makeev}, Sergei Liamaev}
\publication{RecSys 2025}{\href{https://neuralsrg.github.io/publication/2025-08-09-challenge}{Blending Sequential Embeddings, Graphs, and Engineered Features: 4th Place Solution in RecSys Challenge 2025}}{\textbf{Sergei Makeev}, Alexandr Andreev, Vladimir Baikalov, Vladislav Tytskiy, Aleksei Krasilnikov, Kirill Khrylchenko}
\publication{}{\href{https://neuralsrg.github.io/publication/2025-07-12-correcting}{Correcting the LogQ Correction: Revisiting Sampled Softmax for Large-Scale Retrieval}}{Kirill Khrylchenko, Vladimir Baikalov, \textbf{Sergei Makeev}, Artem Matveev, Sergei Liamaev}

\cvsection{Professional Service}
\service{CIKM 2026}{PC member, Demo Track.}
\service{UMAP 2026}{PC member, Industry Track.}
\service{UMAP 2025}{PC member, Industry \& Start-up Track.}

\cvsection{Patents}
\daterow{Dec 2025}{Application for invention \href{https://www1.fips.ru/registers-doc-view/fips_servlet?DB=RUPAT\&DocNumber=2025136959\&TypeFile=html}{2025136959}, Federal Service
for Intellectual Property}{Yandex filed a patent application for Argus (see our paper \href{https://neuralsrg.github.io/publication/2026-07-05-argus}{Scaling Recommender Transformers to One Billion Parameters}).}

\cvsection{Competitions \& Olympiads}
\daterow{Jun 2025}{RecSys Challenge 2025 (RecSys'25)}{Led team \texttt{ambitious} to \href{https://recsys.synerise.com/results/final-leaderboard}{4th place}.}
\daterow{Apr 2024}{Moscow State University, \href{https://confhub.ru/eng/event/7925/}{The Lomonosov Universiade} in Mathematics and Computer Science}{2nd-degree diploma.}
\daterow{Apr 2023}{Moscow State University, \href{https://confhub.ru/eng/event/7925/}{The Lomonosov Universiade} in Mathematics and Computer Science}{1st-degree diploma.}

\cvsection{Talks}
\daterow{Sep 2025}{\href{https://neuralsrg.github.io/talks/2025-09-22-challenge}{Talk at RecSys'25}, O2 Universum, Prague, Czech Republic}{Presented my team's solution to the RecSys Challenge 2025 as an invited speaker.}
\daterow{Mar 2026}{\href{https://neuralsrg.github.io/talks/2026-03-24-challenge_hse}{Talk at HSE University}, Moscow, Russia}{Gave a guest lecture for the \href{https://github.com/KhrylchenkoKirill/DeepRecSys}{DeepRecSys course} taught by Kirill Khrylchenko at the Faculty of Computer Science, HSE University.}

\cvsection{Teaching}
\begin{itemize}
  \compactitem{Contributed to the \href{https://github.com/KhrylchenkoKirill/DeepRecSys}{DeepRecSys undergraduate course} taught by Kirill Khrylchenko at the Faculty of Computer Science, HSE University.}
\end{itemize}

\cvsection{Research in News and Social Media}
\begin{itemize}
  \compactitem{\href{https://neuralsrg.github.io/publication/2026-07-05-argus}{Argus} received broad attention after its launch. It was featured in news outlets (e.g., \href{https://ict.moscow/news/iandeks-vnedriaet-v-rekomendatelnye-sistemy-svoikh-servisov-reshenie-argus-na-baze-generativnogo-ii/}{ict.moscow}, \href{https://yagla.ru/blog/kontekstnaya-reklama/yandeks-predstavil-krupneyshee-obnovlenie-rekomendatelnyh-sistem-ii-budet-zapominat-deystviya-i-interesy-polzovateley-za-god--26/}{yagla.ru}, \href{https://mltimes.ai/yandeks-vnedryaet-novuyu-rekomendatelnuyu-ii-sistemu-argus/}{mltimes.ai}, \href{https://setka.ru/posts/01971ff6-c0ec-4ef4-b81e-a90e7e6a20ba}{setka.ru}, \href{https://toolarium.ru/yandex-argus-rekomendacii/}{toolarium.ru}, \href{https://m.seonews.ru/events/yandeks-zapuskaet-novye-rekomendatelnye-sistemy-na-osnove-generativnykh-modeley/}{m.seonews.ru}, \href{https://rozetked.me/news/39891-yandeks-razrabotal-i-vnedril-novoe-pokolenie-rekomendaciy-na-baze-generativnyh-modeley}{rozetked.me}, \href{https://pikabu.ru/story/yandeks_apnul_rekomendatsii_s_pomoshchyu_generativnyikh_ii_12864778}{pikabu.ru}), covered on \href{https://habr.com/ru/news/920204/}{habr.com}, and discussed on both \href{https://x.com/yandex/status/1935977821622481330}{x.com} and \href{https://habr.com/ru/companies/yandex/articles/919058/}{habr.com}. It was later adapted by our colleagues to many other Yandex products. Its deployment in Yandex Ads also drew extensive discussion (e.g., \href{https://timeweb.com/ru/community/articles/yandeks-vnedril-rekomendatelnuyu-tehnologiyu-argus-v-reklamnye-sistemy}{timeweb.com}, \href{https://habr.com/ru/news/1024284/}{habr.com}, \href{https://adindex.ru/news/digital/2026/04/14/344192.phtml}{adindex.ru}, \href{https://habr.com/ru/companies/click/news/1024506/}{habr.com}, \href{https://dyvannyi-it.ru/yandex-argus/}{dyvannyi-it.ru}, \href{https://marketplacegu.ru/contentzavod/blog/view/yandeks-uvelichil-effektivnost-reklamy-na-38-s}{marketplacegu.ru}, \href{https://www.sostav.ru/publication/yandeks-vnedril-v-reklamnuyu-sistemu-rekomendatelnuyu-tekhnologiyu-argus-83114.html}{sostav.ru}, etc.). \href{https://x.com/_reachsumit?s=21}{Sumit Kumar}, author of a popular information retrieval blog (4.1k subscribers on x.com as of Jul 2026), wrote about Argus on \href{https://x.com/_reachsumit/status/1947864588541935620?s=20}{x.com} and listed it in \href{https://open.substack.com/pub/recsys/p/extending-graphrag-to-millions-of?utm_campaign=post-expanded-share&utm_medium=web}{Top Information Retrieval Papers of the Week}.}
  \compactitem{Our paper \href{https://neuralsrg.github.io/publication/2025-07-12-correcting}{Correcting the LogQ Correction: Revisiting Sampled Softmax for Large-Scale Retrieval} also received extensive discussion: \href{https://vc.ru/id15892/2208230-yandeks-predlagaet-novuyu-formulu-logq-korrektsii}{vc.ru}, \href{https://ict.moscow/news/v-iandekse-razrabotali-novyi-metod-obucheniia-ii-modelei-dlia-rekomendatelnykh-sistem/}{ict.moscow}, \href{https://www.kommersant.ru/doc/8023220}{kommersant.ru}, \href{https://naked-science.ru/article/hi-tech/issledovateli-yandeksa}{naked-science.ru}, \href{https://www.cnews.ru/news/line/2025-09-09_issledovateli_yandeksa}{cnews.ru}, \href{https://itsnew.ru/stati/jandeks-predstavil-obnovlennye-funkcii-dokumentov-i-zapustil-mezhdunarodnuyu-olimpiadu-po-ii.html}{itsnew.ru}, etc. It was listed among \href{https://open.substack.com/pub/recsys/p/towards-agentic-rag-with-deep-reasoning?utm_campaign=post-expanded-share&utm_medium=web}{Top Information Retrieval Papers of the Week}.}
  \compactitem{\href{https://neuralsrg.github.io/publication/2026-06-08-gbla}{Gated Bidirectional Linear Attention for Generative Retrieval}, which I co-authored, was also noted by the community on \href{https://www.linkedin.com/posts/singhsidhukuldeep_scaling-recommendation-systems-how-gated-share-7475730774692933634-NpGL/?utm_source=share&utm_medium=member_desktop&rcm=ACoAAF0dG4wBeGuinWR7V6yIHYXOnEmcjRJY_zY}{linkedin.com} and \href{https://x.com/_reachsumit/status/2063843176553783567?s=20}{x.com}. It was listed in Sumit Kumar's popular blog \href{https://open.substack.com/pub/recsys/p/retrieving-interaction-spaces-for?utm_campaign=post-expanded-share&utm_medium=web}{Top Information Retrieval Papers of the Week}.}
  \compactitem{\href{https://neuralsrg.github.io/publication/2026-07-01-long-term-optimization}{Long-Term Optimization for Large-Scale Generative Retrieval with Off-Policy REINFORCE} was also noted by Sumit Kumar on \href{https://x.com/_reachsumit/status/2074365460016140529?s=46}{x.com}.}
\end{itemize}

\cvsection{Mentorship}
\begin{itemize}
  \compactitem{Mentored two interns at Yandex, Oleg Sorokin and Vladislav Dodonov. Both are now core developers on the RecSys R\&D team (see, for example, \href{https://arxiv.org/abs/2606.08604}{Gryphon}).}
\end{itemize}

\cvsection{Grants}
\begin{itemize}
  \compactitem{Worked on cognitive dynamics analysis and internal speech recognition under Russian Science Foundation grants \href{https://rscf.ru/project/LXjuaL8yl_MKyMwV-E0AYmYXKZ7dVATTIy5v_6WfDOq3PZJVv6H2gRUOxla_daUJ4bgTN-xp/}{\#20-18-00067} and \href{https://rscf.ru/project/OwNy03bhBYAiNgcRD1Lq9DYlzpziFaDtdi0LjD-lDfhmGp1fntBfNIA3DdjvFBf32QgMbnrg/}{\#20-18-00067-\(\Pi\)} as a research assistant at Moscow State University, under the supervision of \href{https://www.researchgate.net/profile/Vartanov-Alexander}{Dr. Alexander Vartanov}.}
\end{itemize}

\cvsection{Technical Skills}
\begin{tabularx}{\textwidth}{@{}p{1.18in}X@{}}
  Languages & Current: Python. Past: C/C++, Assembly, CUDA.\\
  Frameworks & Current: PyTorch. Past: TensorFlow, JAX.\\
\end{tabularx}

\end{document}
